\section{中心胞元约化}
\label{sec:reduction}

本节固定 $\beta$。条带 $B$ 在 $x$ 方向无界，但缺陷只占据中心有界部分。本节把导模问题
等价地写成中心胞元上的问题，其中两侧半波导的作用由它们在人工界面上的 Cauchy 数据空间
表示。与 Fliss 的 DtN 约化\cite[定理~4.5]{Fliss2013}相比，这里不要求半波导 Dirichlet
问题唯一可解；DtN 形式是 Cauchy 数据空间为图像时的特例（第~\ref{subsec:dtn-form}~小节）。

\subsection{胞元与人工界面}
\label{subsec:cells}

记 $e_x=(1,0)$。中心胞元及左右第 $j$ 个周期胞元（$j\geq0$）为
\begin{align}
  \mathcal C_0&:=(x_-,x_+)\times I_d,
  \notag\\
  \mathcal C_{j,+}&:=(x_++jL_+,\,x_++(j+1)L_+)\times I_d,
  &\mathcal C_{j,-}&:=(x_--(j+1)L_-,\,x_--jL_-)\times I_d,
  \label{eq:cells}
\end{align}
并简记 $\mathcal C_\pm:=\mathcal C_{0,\pm}$。于是 $\mathcal C_{j,\pm}=\mathcal C_\pm\pm jL_\pm e_x$，
左右半条带为
\begin{align}
  B_\pm
  &:=\operatorname{int}\bigcup_{j\geq0}\overline{\mathcal C_{j,\pm}},
  &B&=\operatorname{int}\bigl(\overline{B_-}\cup\overline{\mathcal C_0}\cup\overline{B_+}\bigr).
  \label{eq:half-guides}
\end{align}
记 $\Omega_0\subset\mathcal C_0$ 为中心介质柱（允许 $\Omega_0=\varnothing$），
$\Omega_\pm\subset\mathcal C_\pm$ 为基本周期胞元中的介质柱，
$\Omega_{j,\pm}:=\Omega_\pm\pm jL_\pm e_x$。按第~\ref{sec:modelproblem}~节的假定，每个介质柱的
闭包含于相应胞元内部。\pending{左右引线不同（2026-11-04 修改）}$\Omega_-$ 与 $\Omega_+$
分别属于左右两块晶体，形状、系数 $\rho_{\pm,\mathrm c}$ 和胞元宽度 $L_\pm$ 都可以不同。

人工界面与法向约定为
\begin{align}
  \Gamma_\pm&:=\{x_\pm\}\times I_d,
  &\nu_\pm&:=\pm e_x,
  &\Gamma_{j,\pm}&:=\{x_\pm\pm(j+1)L_\pm\}\times I_d,\quad j\geq0.
  \label{eq:artificial-interfaces}
\end{align}
$\Gamma_\pm$ 分隔 $\mathcal C_0$ 与 $\mathcal C_\pm$，$\Gamma_{j,\pm}$ 分隔
$\mathcal C_{j,\pm}$ 与 $\mathcal C_{j+1,\pm}$。$\nu_\pm$ 是 $\mathcal C_0$ 在
$\Gamma_\pm$ 上的外法向，指向半条带 $B_\pm$；它与 $B_\pm$ 自身的外法向相反。下文在
$\Gamma_\pm$ 上的法向导数一律沿 $\nu_\pm$ 计算，无论函数定义在哪一侧。

\subsection{局部解空间与 Cauchy 数据}

对 $U\in\{\mathcal C_0,B_-,B_+\}$，令
\begin{align}
  H^1_\beta(U)
  &:=\bigl\{u\in H^1(U):u|_{y=d/2}=\ee^{\ii\beta d}u|_{y=-d/2}\bigr\},
  \notag\\
  H^{1/2}_\beta(\Gamma_\pm)
  &:=\gamma_\pm H^1_\beta(\mathcal C_0)=\gamma_\pm H^1_\beta(B_\pm),
  &H^{-1/2}_\beta(\Gamma_\pm)&:=H^{1/2}_\beta(\Gamma_\pm)^*,
  \label{eq:trace-spaces}
\end{align}
其中 $\gamma_\pm$ 为 $\Gamma_\pm$ 上的 Dirichlet 迹，迹空间取商范数。这样，$y$ 方向的
准周期兼容性已编码在迹空间中。对 $u\in H^1_\beta(U)$ 且 $\Delta u\in L^2(U)$、$\partial_y u$
满足准周期条件，沿 $\nu_\pm$ 的法向迹 $\partial_{\nu_\pm}u|_{\Gamma_\pm}\in
H^{-1/2}_\beta(\Gamma_\pm)$ 由 Green 公式弱定义。记 Cauchy 数据映射
\begin{align}
  \mathscr T_\pm u
  &:=\bigl(\gamma_\pm u,\ \partial_{\nu_\pm}u|_{\Gamma_\pm}\bigr)
  \in H^{1/2}_\beta(\Gamma_\pm)\times H^{-1/2}_\beta(\Gamma_\pm).
  \label{eq:cauchy-map}
\end{align}

对 $k>0$，定义中心胞元与半波导上的局部解空间
\begin{align}
  \mathcal U_0(k)
  &:=\bigl\{u\in H^1_\beta(\mathcal C_0):-\Delta u=k^2\rho u\text{ 于 }\mathcal C_0,\
    \partial_yu\text{ 满足 $\beta$-准周期条件}\bigr\},
  \notag\\
  \mathcal V_\pm(k)
  &:=\bigl\{v\in H^1_\beta(B_\pm):-\Delta v=k^2\rho_\pm v\text{ 于 }B_\pm,\
    \partial_yv\text{ 满足 $\beta$-准周期条件}\bigr\},
  \label{eq:local-solution-spaces}
\end{align}
方程均按分布理解，因此包含了内部材料界面上的传输条件；$\Gamma_\pm$ 上不施加任何条件。
$\mathcal V_\pm(k)$ 中的函数属于 $H^1(B_\pm)$，因而在半无限方向上平方可积。半波导的
\emph{Cauchy 数据空间}定义为
\begin{align}
  \mathcal G_\pm(k)
  &:=\mathscr T_\pm\mathcal V_\pm(k)
  \subset H^{1/2}_\beta(\Gamma_\pm)\times H^{-1/2}_\beta(\Gamma_\pm).
  \label{eq:cauchy-data-space}
\end{align}

\begin{lemma}[Cauchy 数据的唯一延拓]
  \label{lem:unique-continuation}
  若 $v\in\mathcal V_\pm(k)$ 且 $\mathscr T_\pm v=0$，则 $v=0$。
\end{lemma}

\begin{proof}
  以 $+$ 侧为例。由于介质柱不与 $\Gamma_+$ 相交，存在 $\epsilon>0$ 使
  $(x_+-\epsilon,x_++\epsilon)\times I_d$ 中的系数恒为 $\rho_{\mathrm e}$。把 $v$ 在
  $(x_+-\epsilon,x_+)\times I_d$ 上延拓为零，并沿 $y$ 准周期延拓。由于两侧的值迹和
  法向迹都相等，延拓后的函数在 $(x_+-\epsilon,x_++\epsilon)\times\RR$ 中是
  $-\Delta\tilde v=k^2\rho_{\mathrm e}\tilde v$ 的分布解，因而实解析；它在一个开集上为零，
  故在该带中恒为零。$B_+$ 中背景区域是连通的，由解析延拓 $v$ 在整个背景区域上为零。
  于是每个介质柱边界上外侧 Cauchy 数据为零，由传输条件内侧 Cauchy 数据也为零；由柱内
  Green 表示公式，$v$ 在柱内为零。
\end{proof}

\subsection{等价性}

\begin{proposition}[中心约化]
  \label{prop:center-reduction}
  固定 $\beta$ 与 $k>0$，令
  \begin{align}
    \mathcal U_0^{\mathcal G}(k)
    &:=\bigl\{u_0\in\mathcal U_0(k):
      \mathscr T_-u_0\in\mathcal G_-(k),\ \mathscr T_+u_0\in\mathcal G_+(k)\bigr\}.
    \label{eq:matched-center-space}
  \end{align}
  则限制映射 $u\mapsto u|_{\mathcal C_0}$ 是从 $\ker(A(\beta)-k^2)$ 到
  $\mathcal U_0^{\mathcal G}(k)$ 的线性双射。特别地，$k^2$ 是 $A(\beta)$ 的特征值当且
  仅当 $\mathcal U_0^{\mathcal G}(k)\ne\{0\}$，且两者维数相同。
\end{proposition}

\begin{proof}
  \emph{映射有定义。}设 $u\in\ker(A(\beta)-k^2)$。则 $u|_{\mathcal C_0}\in\mathcal U_0(k)$，
  $u|_{B_\pm}\in\mathcal V_\pm(k)$。由于 $u\in H^1(B)$ 且 $\Delta u\in L^2(B)$，在
  $\Gamma_\pm$ 两侧的值迹相同，沿同一法向 $\nu_\pm$ 的法向迹也相同，故
  $\mathscr T_\pm(u|_{\mathcal C_0})=\mathscr T_\pm(u|_{B_\pm})\in\mathcal G_\pm(k)$。

  \emph{单射。}若 $u|_{\mathcal C_0}=0$，则 $\mathscr T_\pm(u|_{B_\pm})=0$，由
  引理~\ref{lem:unique-continuation} 得 $u|_{B_\pm}=0$，故 $u=0$。

  \emph{满射。}设 $u_0\in\mathcal U_0^{\mathcal G}(k)$，取 $v_\pm\in\mathcal V_\pm(k)$
  使 $\mathscr T_\pm v_\pm=\mathscr T_\pm u_0$（由引理~\ref{lem:unique-continuation}，
  $v_\pm$ 唯一）。把 $v_-$、$u_0$、$v_+$ 拼接为 $B$ 上的函数 $u$。由于 $\Gamma_\pm$ 两侧
  值迹相同，$u\in H^1_\beta(B)$；由于法向迹也相同，分布 $\Delta u$ 在 $\Gamma_\pm$ 上
  没有奇异部分，故 $\Delta u=-k^2\rho u\in L^2(B)$。$y$ 方向的导数准周期条件在每一部分
  上成立。因此 $u\in D(A(\beta))$ 且 $A(\beta)u=k^2u$，其限制为 $u_0$。
\end{proof}

命题~\ref{prop:center-reduction} 对任意 $k>0$ 成立，并不依赖谱隙假设。谱隙的作用在于
刻画 $\mathcal V_\pm(k)$ 的结构：当 $k^2\in G_\beta$ 时（$G_\beta$ 是两条引线的公共谱隙），半波导中不存在沿 $x$ 传播的
Bloch 波，有限能量解逐胞元指数衰减。Joly、Li 与 Fliss 用一个谱半径小于一的胞元传播
算子 $\mathcal R_\pm$ 表示这些解，即解在第 $j$ 个胞元上的限制由界面数据
$\mathcal R_\pm^j\varphi$ 经单胞问题得到\cite{JolyLiFliss2006}；该算子的 Jordan 结构见
\cite{HohageSoussi2013}。第~\ref{sec:discretization}~节的胞元传播矩阵
$\mathsf P_\pm$ 正是这一算子在端口 Rayleigh 坐标中的离散对应。

\subsection{DtN 形式}
\label{subsec:dtn-form}

若对每个 $\varphi\in H^{1/2}_\beta(\Gamma_\pm)$，存在唯一的 $v\in\mathcal V_\pm(k)$ 满足
$\gamma_\pm v=\varphi$，则 $\mathcal G_\pm(k)$ 是 DtN 算子
\begin{align}
  \Lambda_\pm(k):H^{1/2}_\beta(\Gamma_\pm)\to H^{-1/2}_\beta(\Gamma_\pm),
  \qquad
  \Lambda_\pm(k)\varphi:=\partial_{\nu_\pm}v|_{\Gamma_\pm}
  \label{eq:continuous-dtn}
\end{align}
的图像，而 $\mathscr T_\pm u_0\in\mathcal G_\pm(k)$ 等价于透明边界条件
$\partial_{\nu_\pm}u_0=\Lambda_\pm(k)\gamma_\pm u_0$。这就是 Fliss 的约化
\cite[定理~4.5]{Fliss2013}；由于 Fliss 沿半条带自身的外法向定义 DtN，她的算子等于
$-\Lambda_\pm(k)$。在谱隙中，半波导 Dirichlet 问题除一个可数的例外集外适定
\cite[定理~4.1]{Fliss2013}；例外值是在 $\Gamma_\pm$ 上施加齐次 Dirichlet 条件后半波导
算子的特征值。在例外值处 $\mathcal G_\pm(k)$ 不再是图像，DtN 约化失效，而
命题~\ref{prop:center-reduction} 仍然成立。Robin-to-Robin 条件是绕开这一问题的另一种
途径\cite{FlissKlindworthSchmidt2015}。本文以 Cauchy 数据形式为主线；DtN 形式只在
离散 Dirichlet 坐标安全时作为等价的降维形式出现。
